Part 1 — Genuine Analysis: Proposal vs. Notebook, Using the Real Outputs
What the outputs actually say (the story your dissertation should tell)
The data. 21,347 households, 47 counties, 534 engineered columns. Missingness is structural, not sloppy: 112 columns >80% missing are survey routing skips; 126 columns are 100% complete. You cleaned 540 → 446 columns and lost zero rows — that is a defensible, auditable cleaning narrative.

The categorical variables (your specific concern). The distribution is meaningful, not decorative:

Tenure: 61.5% owner-occupiers (13,118), 32.5% renters (6,938) — a genuine dual market.
Gender: 31.9% female-headed — a large, policy-relevant minority.
Settlement: 44.3% urban vs 55.7% rural (urban-oversampled vs Kenya's ~30% — your S0.6 census weighting is what neutralises this; justify it exactly as I've now written it).
Income: 4 KNBS brackets with midpoints 10k/32.5k/100k/225k. Critical honest finding: this bracket quantisation is the #1 SHAP driver (0.0517) — partly mechanical (it's HCB's denominator) but partly a data granularity finding you should own: finer income measurement would sharpen targeting.
Hazard: only 1.3% hazard-prone — which is why D3 had to be re-engineered (structural exposure 65.4%, climate service failures 33.5%). This is a strength: you traded a sparse flag for dense behavioural signals.
The categorical variables that matter vs. those that don't. Your SHAP ranking settles it: income (0.0517) ≫ tenure (0.0211) ≫ water spend (0.0100) > floor area (0.0076) > approved building (0.0068); urban/rural is only 13th (0.0030) — its effect is absorbed by tenure/utilities/county context. And the SHAP-derived pillar weights β = (0.241, 0.352, 0.129, 0.179, 0.099) say the data prizes tenure insecurity above everything — a genuinely non-obvious, defensible finding.

The headline distribution. HFVS tiers: Low 10.0% / Moderate 43.0% / High 35.6% / Critical 11.4%. Against the 4.68% flagged by a raw HCB>0.30 screen — a cost-only rule sees 4.68%, your pipeline sees 47%. That sentence is your dissertation's thesis in one line.

Model chain (real numbers). Linear 0.2285 → LightGBM 0.5923 → ensemble 0.6020 → hurdle champion 0.6172 (Δ+0.389). Quantile band 0.0352 (tight). DEA: mean θ_CCR 0.832, 9/47 on the frontier. MILP: 6,000 units, KSh 19.75B — and the two regimes converge at identical cost, while at the official 200,000-unit target the whole 15,365-unit backlog is fundable (KSh 50.58B of KSh 73.2B) → delivery capacity (1,795 units actually completed FY24/25), not levy revenue, is the binding constraint. That is your most publishable policy insight and it was buried in a code printout — it's now in the conclusions markdown.

Have you left anything out of the previous feedback? (Before my edits)
Feedback item	Status
D3 sparse-hazard fix (02_D3_FIX.md)	✅ done
Leakage gates (06_LEAKAGE.md)	✅ done & asserted
Eq.2 SHAP-derived weights promise-vs-delivery gap	✅ done (round 3)
Residual sign labels swapped	✅ fixed (round 3)
Hardcoded stale R² in summary table	✅ fixed (round 2)
S2.1 placeholder "X%"	❌ was left — fixed now (31.9%)
S9.5 limitations table contradicted the code (said weights equal, cost placeholder)	❌ was left — fixed now
S7 falsely claimed the Q1 correction "is applied in S8"	❌ was left — fixed now
S8 observation claimed wide Critical-tier bands (contradicted: 0.0% above 0.20 width)	❌ was left — fixed now
S11 "84% near-zero burden" unverified	❌ — replaced with the real 61.5% owner + median-vs-mean evidence
What I changed in the build scripts (all edits are markdown chunks; rebuild verified, 93 cells)
S0: Executive framing (Business Question / Decision to be Made) written into the Problem Statement, anchored on real values (mean HCB 0.090 vs median 0.0133, 4.68% stressed vs 47% High+Critical). Headline-results row added to the title card.
S0.6: Explicit "why we use population data" — converts household shares into 2026 people, protects against the urban oversample, makes MILP umax caps defensible.
S2 (categorical landscape): every observation now carries real values (31.9% female-headed; 44.3/55.7 split; tenure means/medians 0.1032/0.0178 vs 0.0732/0.0089; 1.3% hazard-prone; 15 age outliers).
S3/S4: shape-journey narrated with printed shapes; S4.2 rewritten as the honest "families are NOT redundant" result; new S4.3 narration cell for the SHAP weights, discriminant validity (max r=0.469) and tier distribution.
S6/S7/S8: SHAP rank-order story with real values; Q1 inequity quantified (+0.0298, MAE 0.2020) and the correction reframed as a documented policy choice, not applied; tuning failure and tight quantile fan reframed as defensible honest results.
S9.5: RO table now carries the real outcome numbers; limitations table updated (two items marked RESOLVED, three new limitations: regime convergence, median-imputation, shapefile fallback).
S10/S11: goals restated against the real 0.5923 champion; hurdle architecture justified by the structural-zero evidence and its 0.6172/0.6075 results.
Part 2 — How to Reach HIGH MERIT (depth, practicability, applicability, product)
Run the two declared sensitivity re-runs before defence — they're promised in your own markdown: (a) MILP at KSh 1.05M–1.95M/unit with a county-rank Spearman; (b) a budget level where Regime A ≠ Regime B. Both are cheap; both close the biggest remaining "so what" gaps.
Get the county shapefile (KNBS/gadm-level) — one choropleth replaces two bar charts and materially lifts the "product" criterion. The code already falls back gracefully.
Reframe MILP's finding as the centrepiece: "revenue is not the constraint; delivery capacity is (1,795 completions vs a KSh 73.2B levy)" — connect to the DEA result (9 frontier counties = where capacity exists today). This makes Tier 2 and Tier 3 talk to each other, which is what depth looks like.
Own the Q1 bias as a design choice: compare the hurdle model's Q1 residual explicitly — if the hurdle reduces Q1 error, that's a result; report it either way.
Product: freeze MODEL_RESULTS into a one-page model card (champion, features, bias audit, refresh cadence) and reference it from the proposal — examiners reward deployability evidence.
Part 3 — Executive Communication Model (real values, in order)
Business Question. Which Kenyan households and counties are most financially vulnerable in housing, and how should the Housing Levy's next tranche be allocated to buy the greatest vulnerability reduction per shilling?
Decision to be Made. County-level allocation of housing units and budget (where to build, how much, under what subsidy mix) — replacing a manual, cost-centric workflow.
What the Data is Telling Us. 21,347 households: mean HCB 0.090 but median 0.0133; only 4.68% cross the 30% stress line, yet 47% are High/Critical on the composite. Renters (32.5%) carry higher insecurity; 31.9% of households are female-headed; hazard exposure is sparse (1.3%) but service failures are common (33.5%).
Model Insight. Tenure insecurity is the heaviest HFVS pillar (β₂=0.352). HCB prediction champion: two-stage hurdle, R²=0.6172 (+0.389 over linear); drivers: income bracket (0.0517), tenure (0.0211), water spend (0.0100). County DEA: mean θ_CCR 0.832; 9 frontier counties (Bungoma, Garissa, Kiambu, Lamu, Migori, Nairobi, Nyeri, Tharaka-Nithi, Wajir); weakest: West Pokot (0.550), Nyamira, Marsabit.
Confidence / Uncertainty. Test-set metrics on a leakage-safe split; prediction bands tight (mean 0.0352 Q75–Q25; 0% of households with Q90–Q10 > 0.20). The material uncertainty is equity, not variance: Q1 households underpredicted (mean residual +0.0298, MAE 0.2020).
Business Impact. Delivery-constrained MILP: 6,000 units across 47 counties for KSh 19.75B — only 27% of the KSh 73.2B levy; at the official 200,000-unit target the entire modelled backlog (15,365 units, KSh 50.58B) is fundable. Every efficiency point recovered in the 38 sub-frontier counties converts directly into unspent levy or extra units.
Risks & Limitations. Q1 underprediction; county predictors median-imputed (land prices: 7,315 nulls); single survey wave (no trend); no shapefile (bar-chart maps this run); MILP regimes converge at this quota; construction cost sensitivity band not yet re-run.
Management Options. (a) Concentrate on efficient counties; (b) universal coverage — identical cost at this quota; (c) target High-HFVS/low-θ counties with dual capacity + voucher interventions; (d) fix delivery capacity, not revenue.
Recommended Action. Universal coverage at the delivery-constrained quota (it is free), with the 30% voucher ring-fence and capacity-building concentrated in the High-HFVS/low-efficiency quadrant, and Q1 predictions upward-adjusted via the documented correction.
Monitor Results. Recalibrate HFVS annually with each KHS wave; re-estimate DEA θ yearly as permit/EIA data updates; track completions vs the 1,795 baseline; re-run the cost sensitivity band each construction-price cycle.
Bottom line: the notebook now tells the truth with the real numbers and reads as an executive product. 